\documentclass[10pt,a4paper]{amsart}
\usepackage{amsmath,amssymb,amsthm,graphicx}
\usepackage{verbatim}
\newenvironment{reference}{\comment}{\endcomment}
%\newenvironment{reference}{}{}
\newenvironment{slogan}{\comment}{\endcomment}
\newenvironment{history}{\comment}{\endcomment}

% For commutative diagrams we use Xy-pic
\usepackage[all]{xy}

% We use 2cell for 2-commutative diagrams.
\xyoption{2cell}
% Theorem environments.
%
\theoremstyle{plain}
\newtheorem{theorem}[subsection]{Theorem}
\newtheorem{proposition}[subsection]{Proposition}
\newtheorem{lemma}[subsection]{Lemma}

\theoremstyle{definition}
\newtheorem{definition}[subsection]{Definition}
\newtheorem{example}[subsection]{Example}
\newtheorem{exercise}[subsection]{Exercise}
\newtheorem{situation}[subsection]{Situation}

\theoremstyle{remark}
\newtheorem{remark}[subsection]{Remark}
\newtheorem{remarks}[subsection]{Remarks}

\numberwithin{equation}{subsection}

% Macros
%
\def\lim{\mathop{\mathrm{lim}}\nolimits}
\def\colim{\mathop{\mathrm{colim}}\nolimits}
\def\Spec{\mathop{\mathrm{Spec}}}
\def\Hom{\mathop{\mathrm{Hom}}\nolimits}
\def\Ext{\mathop{\mathrm{Ext}}\nolimits}
\def\SheafHom{\mathop{\mathcal{H}\!\mathit{om}}\nolimits}
\def\SheafExt{\mathop{\mathcal{E}\!\mathit{xt}}\nolimits}
\def\Sch{\mathit{Sch}}
\def\Mor{\mathop{\mathrm{Mor}}\nolimits}
\def\Ob{\mathop{\mathrm{Ob}}\nolimits}
\def\Sh{\mathop{\mathit{Sh}}\nolimits}
\def\NL{\mathop{N\!L}\nolimits}
\def\CH{\mathop{\mathrm{CH}}\nolimits}
\def\proetale{{pro\text{-}\acute{e}tale}}
\def\etale{{\acute{e}tale}}
\def\QCoh{\mathit{QCoh}}
\def\Ker{\mathop{\mathrm{Ker}}}
\def\Im{\mathop{\mathrm{Im}}}
\def\Coker{\mathop{\mathrm{Coker}}}
\def\Coim{\mathop{\mathrm{Coim}}}

% Boxtimes
%
\DeclareMathSymbol{\boxtimes}{\mathbin}{AMSa}{"02}

%
% Macros for moduli stacks/spaces
%
\def\QCohstack{\mathcal{QC}\!\mathit{oh}}
\def\Cohstack{\mathcal{C}\!\mathit{oh}}
\def\Spacesstack{\mathcal{S}\!\mathit{paces}}
\def\Quotfunctor{\mathrm{Quot}}
\def\Hilbfunctor{\mathrm{Hilb}}
\def\Curvesstack{\mathcal{C}\!\mathit{urves}}
\def\Polarizedstack{\mathcal{P}\!\mathit{olarized}}
\def\Complexesstack{\mathcal{C}\!\mathit{omplexes}}
% \Pic is the operator that assigns to X its picard group, usage \Pic(X)
% \Picardstack_{X/B} denotes the Picard stack of X over B
% \Picardfunctor_{X/B} denotes the Picard functor of X over B
\def\Pic{\mathop{\mathrm{Pic}}\nolimits}
\def\Picardstack{\mathcal{P}\!\mathit{ic}}
\def\Picardfunctor{\mathrm{Pic}}
\def\Deformationcategory{\mathcal{D}\!\mathit{ef}}
\usepackage{cleveref}
\crefname{originallemma}{original lemma}{original lemmas}
\Crefname{originallemma}{Original lemma}{Original lemmas}
\makeatletter
\newlabel{algebra-lemma-standard-open-containing-maximal-point}{{00EV}{source}{}{sourcepointer-1}{}}
\newlabel{algebra-lemma-standard-open-containing-maximal-point@cref}{{[originallemma][0][]00EV}{[1][0][]source}{}{}{}}
\makeatother
\crefname{originallemma}{original lemma}{original lemmas}
\Crefname{originallemma}{Original lemma}{Original lemmas}
\makeatletter
\newlabel{decent-spaces-lemma-decent-generic-points}{{0ABV}{source}{}{sourcepointer-2}{}}
\newlabel{decent-spaces-lemma-decent-generic-points@cref}{{[originallemma][0][]0ABV}{[1][0][]source}{}{}{}}
\makeatother
\crefname{originallemma}{original lemma}{original lemmas}
\Crefname{originallemma}{Original lemma}{Original lemmas}
\makeatletter
\newlabel{decent-spaces-lemma-filter-quasi-compact-quasi-separated}{{07ST}{source}{}{sourcepointer-3}{}}
\newlabel{decent-spaces-lemma-filter-quasi-compact-quasi-separated@cref}{{[originallemma][0][]07ST}{[1][0][]source}{}{}{}}
\makeatother
\crefname{originalsection}{original section}{original sections}
\Crefname{originalsection}{Original section}{Original sections}
\makeatletter
\newlabel{decent-spaces-section-reasonable-decent}{{03I7}{source}{}{sourcepointer-4}{}}
\newlabel{decent-spaces-section-reasonable-decent@cref}{{[originalsection][0][]03I7}{[1][0][]source}{}{}{}}
\makeatother
\crefname{originallemma}{original lemma}{original lemmas}
\Crefname{originallemma}{Original lemma}{Original lemmas}
\makeatletter
\newlabel{properties-lemma-generic-point-in-constructible}{{0AAW}{source}{}{sourcepointer-5}{}}
\newlabel{properties-lemma-generic-point-in-constructible@cref}{{[originallemma][0][]0AAW}{[1][0][]source}{}{}{}}
\makeatother
\crefname{originaldefinition}{original definition}{original definitions}
\Crefname{originaldefinition}{Original definition}{Original definitions}
\makeatletter
\newlabel{spaces-properties-definition-dimension-local-ring}{{04NA}{source}{}{sourcepointer-6}{}}
\newlabel{spaces-properties-definition-dimension-local-ring@cref}{{[originaldefinition][0][]04NA}{[1][0][]source}{}{}{}}
\makeatother
\crefname{originallemma}{original lemma}{original lemmas}
\Crefname{originallemma}{Original lemma}{Original lemmas}
\makeatletter
\newlabel{spaces-properties-lemma-codimension-0-points}{{0BAQ}{source}{}{sourcepointer-7}{}}
\newlabel{spaces-properties-lemma-codimension-0-points@cref}{{[originallemma][0][]0BAQ}{[1][0][]source}{}{}{}}
\makeatother
\crefname{originalproposition}{original proposition}{original propositions}
\Crefname{originalproposition}{Original proposition}{Original propositions}
\makeatletter
\newlabel{spaces-properties-proposition-locally-quasi-separated-open-dense-scheme}{{06NH}{source}{}{sourcepointer-8}{}}
\newlabel{spaces-properties-proposition-locally-quasi-separated-open-dense-scheme@cref}{{[originalproposition][0][]06NH}{[1][0][]source}{}{}{}}
\makeatother
\title{Characterizations of points of codimension zero on a quasi-separated algebraic space}
\author{The Stacks Project Authors}
\date{stacks-project:spaces-over-fields.tex:local-snapshot-20261012; GNU FDL 1.2 or later}
\begin{document}
\maketitle
\noindent\textit{Editorial note.} All delivered mathematical text is the author's own, copied from the pinned Stacks Project chapter file spaces-over-fields.tex; no statement or proof step was rewritten, reordered or removed.\par\medskip
\noindent\textit{Editorial note.} The counted proof is the author's complete argument: the identification of the first two conditions by definition, the elementary comparison with generic points through the filtration of a quasi-separated decent space and the standard open neighbourhood on which the chart is an isomorphism, and the final appeal to the proposition on dense open subschemes.\par\medskip
\noindent\textit{Editorial note.} The label lemma-generic-point-in-schematic-locus names the counted statement, and the proof environment immediately following it in the source is delivered verbatim as its proof.\par\medskip
\noindent\textit{Editorial note.} The counted statement is carried by the author's own lemma environment in the chapter source.\par\medskip
\noindent\textit{Editorial note.} Pre-open gate: python operations/agent-authoring/tools/precard\_gate.py --source spaces-over-fields.tex --label lemma-generic-point-in-schematic-locus --start 1597 --end 1653 --local-id 28057 returned RESULT=OK (L28057 free, occupancy FREE, public interval clear). Candidate taken from sources/stacks-spaces-over-fields/candidates-halfA.json.\par\medskip
\noindent\textit{Editorial note.} The proof notes that a more conceptual proof through decent spaces is available and then supplies 'a more elementary proof of the equivalence', which is delivered in full, so the statement is counted.\par\medskip
\noindent\textit{Editorial note.} The statement and proof are copied from the pinned chapter file; the cited inputs are kept as external pointers with their canonical Stacks tags.\par\medskip
\noindent\textit{Allowed standard prerequisites (editorial).}\begin{itemize}
\item Concise description of a quasi-separated decent space. Decent Spaces, the lemma providing a filtration by open subspaces with bijective maps from quasi-compact schemes used in the elementary proof. Supplies the chain of open subspaces and the maps f\_i whose points compute the given point.
\item Generic points of schemes and constructible subsets. Properties, the lemma that a generic point of a scheme has an open neighbourhood disjoint from a closed subset not containing it, and Same chapter, the lemma describing codimension zero points of a scheme as generic points of irreducible components. Supply the open neighbourhood W on which the map f\_i is an isomorphism and the identification of codimension zero points with generic points.
\item The open dense subscheme of a locally quasi-separated space. Properties of Spaces, the proposition producing a dense open subspace which is a scheme containing every point of codimension zero. Supplies the final statement of the lemma.
\end{itemize}\medskip
\ifdefined\hypertarget\hypertarget{sourcepointer-1}{}\fi
\noindent\textit{Source reference.} Original lemma 00EV: sources/stacks-chow/tags.txt, source lines 549--549. Prerequisite pointer invoked at proof line(s) 1640; it names a standard fact of a sibling chapter that the counted argument uses, and every step of the counted proof built on it is retained verbatim in the delivered body.\par\medskip
\ifdefined\hypertarget\hypertarget{sourcepointer-2}{}\fi
\noindent\textit{Source reference.} Original lemma 0ABV: sources/stacks-chow/tags.txt, source lines 13033--13033. Prerequisite pointer invoked at proof line(s) 1612; it names a standard fact of a sibling chapter that the counted argument uses, and every step of the counted proof built on it is retained verbatim in the delivered body.\par\medskip
\ifdefined\hypertarget\hypertarget{sourcepointer-3}{}\fi
\noindent\textit{Source reference.} Original lemma 07ST: sources/stacks-chow/tags.txt, source lines 9913--9913. Prerequisite pointer invoked at proof line(s) 1629; it names a standard fact of a sibling chapter that the counted argument uses, and every step of the counted proof built on it is retained verbatim in the delivered body.\par\medskip
\ifdefined\hypertarget\hypertarget{sourcepointer-4}{}\fi
\noindent\textit{Source reference.} Original section 03I7: sources/stacks-chow/tags.txt, source lines 4593--4593. Prerequisite pointer invoked at proof line(s) 1614; it names a standard fact of a sibling chapter that the counted argument uses, and every step of the counted proof built on it is retained verbatim in the delivered body.\par\medskip
\ifdefined\hypertarget\hypertarget{sourcepointer-5}{}\fi
\noindent\textit{Source reference.} Original lemma 0AAW: sources/stacks-chow/tags.txt, source lines 12999--12999. Prerequisite pointer invoked at proof line(s) 1638; it names a standard fact of a sibling chapter that the counted argument uses, and every step of the counted proof built on it is retained verbatim in the delivered body.\par\medskip
\ifdefined\hypertarget\hypertarget{sourcepointer-6}{}\fi
\noindent\textit{Source reference.} Original definition 04NA: sources/stacks-chow/tags.txt, source lines 6071--6071. Prerequisite pointer invoked at proof line(s) 1621; it names a standard fact of a sibling chapter that the counted argument uses, and every step of the counted proof built on it is retained verbatim in the delivered body.\par\medskip
\ifdefined\hypertarget\hypertarget{sourcepointer-7}{}\fi
\noindent\textit{Source reference.} Original lemma 0BAQ: sources/stacks-chow/tags.txt, source lines 14218--14218. Prerequisite pointer invoked at proof line(s) 1633; it names a standard fact of a sibling chapter that the counted argument uses, and every step of the counted proof built on it is retained verbatim in the delivered body.\par\medskip
\ifdefined\hypertarget\hypertarget{sourcepointer-8}{}\fi
\noindent\textit{Source reference.} Original proposition 06NH: sources/stacks-chow/tags.txt, source lines 8536--8536. Prerequisite pointer invoked at proof line(s) 1652; it names a standard fact of a sibling chapter that the counted argument uses, and every step of the counted proof built on it is retained verbatim in the delivered body.\par\medskip
\begin{lemma}
\label{lemma-generic-point-in-schematic-locus}
Let $S$ be a scheme. Let $X$ be a quasi-separated algebraic space over $S$.
Let $x \in |X|$. The following are equivalent
\begin{enumerate}
\item $x$ is a point of codimension $0$ on $X$,
\item the local ring of $X$ at $x$ has dimension $0$, and
\item $x$ is a generic point of an irreducible component of $|X|$.
\end{enumerate}
If true, then there exists an open subspace of $X$
containing $x$ which is a scheme.
\end{lemma}

\begin{proof}
The equivalence of (1), (2), and (3) follows from
Decent Spaces, Lemma \ref{decent-spaces-lemma-decent-generic-points}
and the fact that a quasi-separated algebraic space is decent
(Decent Spaces, Section \ref{decent-spaces-section-reasonable-decent}).
However in the next paragraph we will give a more elementary proof of the
equivalence.

\medskip\noindent
Note that (1) and (2) are equivalent by definition
(Properties of Spaces, Definition
\ref{spaces-properties-definition-dimension-local-ring}).
To prove the equivalence of (1) and (3) we may assume $X$ is quasi-compact.
Choose
$$
\emptyset = U_{n + 1} \subset
U_n \subset U_{n - 1} \subset \ldots \subset U_1 = X
$$
and $f_i : V_i \to U_i$ as in Decent Spaces, Lemma
\ref{decent-spaces-lemma-filter-quasi-compact-quasi-separated}.
Say $x \in U_i$, $x \not \in U_{i + 1}$. Then $x = f_i(y)$ for
a unique $y \in V_i$. If (1) holds, then $y$ is a generic point of
an irreducible component of $V_i$ (Properties of Spaces, Lemma
\ref{spaces-properties-lemma-codimension-0-points}).
Since $f_i^{-1}(U_{i + 1})$ is a quasi-compact open of $V_i$
not containing $y$, there is an open neighbourhood $W \subset V_i$
of $y$ disjoint from $f_i^{-1}(V_i)$
(see
Properties, Lemma \ref{properties-lemma-generic-point-in-constructible}
or more simply Algebra, Lemma
\ref{algebra-lemma-standard-open-containing-maximal-point}).
Then $f_i|_W : W \to X$ is an isomorphism onto its image and hence
$x = f_i(y)$ is a generic point of $|X|$. Conversely, assume (3) holds.
Then $f_i$ maps $\overline{\{y\}}$ onto the irreducible component
$\overline{\{x\}}$ of $|U_i|$. Since $|f_i|$ is bijective over
$\overline{\{x\}}$, it follows that $\overline{\{y\}}$
is an irreducible component of $U_i$. Thus $x$ is a point of
codimension $0$.

\medskip\noindent
The final statement of the lemma is
Properties of Spaces, Proposition
\ref{spaces-properties-proposition-locally-quasi-separated-open-dense-scheme}.
\end{proof}

\end{document}
